\subsection*{40. Researchers Are Paid to Reduce Uncertainty, Not to Produce Impressive Conclusions}

A researcher may create a new theorem, evaluate a scientific hypothesis, develop a model, design a rigorous experiment, or turn conflicting findings into a defensible conclusion. The deepest shared competence is \textbf{epistemic discipline}: distinguishing evidence from interpretation, identifying assumptions, entertaining rival explanations, and changing views in response to results.

A research hiring process should therefore examine problem selection, technical depth, methods, data integrity, reproducibility, communication, and judgment about when \textbf{not} to make a claim. Publications and citations can offer context but are field- and opportunity-dependent; authorship contribution may be hard to infer. Ask about a result that failed, a model the candidate abandoned, and a conclusion they revised. Strong researchers can describe their uncertainty without dissolving into indecision.

\subsection*{41. Use a Research Critique Rather Than a Research Trivia Contest}

Provide an anonymized 2-page study brief: promising effect size, small sample, unclear exclusions, a possible leakage path, and no preregistered analysis. Ask the candidate for the most important concerns, what further data would resolve them, and whether a public claim is justified. For theoretical work, give a flawed proof sketch and ask which lemma fails. For computational research, request a reproducibility plan that states seeds, environments, baselines, evaluation metrics, and robustness checks.

Score the ability to prioritize substantive validity threats. A candidate who names every possible statistical objection without judgment may be less useful than 1 who identifies the 2 decisive assumptions and designs an efficient test. The work sample should match the research domain: laboratory safety, field logistics, or theorem proving cannot be collapsed into the same task.

\subsection*{42. Distinguish Research Originality from Performance on Familiar Problems}

Many candidates can reproduce fashionable methods or extend an obvious benchmark. Originality includes framing a consequential question, noticing what others overlook, proposing a nontrivial mechanism, and designing tests whose failure would be informative. It is difficult to measure in a short interview. Ask candidates to select 1 open problem in a bounded domain, explain why it matters, summarize the strongest alternatives, and propose a 1st falsifiable investigation under a fixed budget.

Do not demand that applicants disclose unpublished confidential work or donate a proprietary research agenda. Use fictional datasets, public literature, or compensated exercises. Senior research hiring should include independent technical referees who understand the domain and can evaluate depth. Committee prestige is no substitute for the ability to engage the substance of a candidate's work.

\subsection*{43. A Research Scorecard}

For a research scientist in applied statistics or AI, an illustrative scorecard is:
\begin{itemize}[leftmargin=*]
\item Methodological validity: 25\%
\item Technical depth and independence: 20\%
\item Problem choice and formulation: 20\%
\item Reproducibility and documentation: 15\%
\item Critical self-correction: 10\%
\item Communication and collaboration: 10\%
\end{itemize}
Require a non-negotiable pass on research integrity, data stewardship, and fundamental field competence.

The weighting should change with the job. A research engineer may need much more software and infrastructure competence; a foundational theoretician may need proof technique; an applied industrial researcher may need client problem translation and useful deployment. Include evidence from research artifacts, peer critique, a narrow technical exercise, references, and discussion of failures. Avoid treating publication count as a score that can compensate for untrustworthy methods.

\subsection*{44. Hire Science Talent for Questions the Institution Can Responsibly Support}

A young venture may be attracted to a distinguished researcher even though it lacks time, data rights, lab infrastructure, peers, grant support, or a coherent research direction. That is not necessarily an opportunity; it may be a way to spend money on a prestigious title without creating a research enterprise. Selection must be accompanied by an honest account of the environment that will exist.

Research careers depend on intellectual freedom, mentoring, recognition, ethical review, equipment, long-term continuity, and a credible path to dissemination. If those conditions do not exist, hire bounded collaborators or advisers for well-defined work rather than promising an institutional research program prematurely. Eventually a research-active Bayes Institute must learn to select for scientific rigor \textbf{and} build the system in which rigorous science can be produced.
